%!TEX root = slides.tex

\title{Data Representation and Querying}
\subtitle{}
\author{ian.mcloughlin@gmit.ie}
\date{}


\begin{frame}
	\titlepage
\end{frame}

\begin{frame}
	\frametitle{Topics}
	\tableofcontents
\end{frame}

\section{About this module}


\begin{frame}{Learning outcomes}
	On completion of this module the learner will/should be able to:
	\begin{itemize}
		\item Explain the benefits and the limitations of a variety of data models.
		\item Determine the most appropriate data model given a set of requirements.
		\item Represent, integrate and querying large datasets using existing API's and frameworks.
		\item Describe the principles behind both the linked data and the open data movements.
	\end{itemize}
\end{frame}


\begin{frame}{Examinations}
	\begin{table}
		\begin{tabular}{p{4cm}r@{\hspace{0.5cm}}p{4cm}}
			Type & \% & Date \\
			\hline
			Project & 50 & Week 8 \\
			End of Semester Exam & 50 & See exams timetable
		\end{tabular}
	\end{table}
\end{frame}


\section{HTTP}

\begin{frame}{HyperText Transfer Protocol}
  \begin{description}
		\item[HyperText] Text with links.
    \vspace{0.25cm}
		\item[Transfer] Communication of data.
    \vspace{0.25cm}
		\item[Protocol] Set of rules for communication.
  \end{description}
\end{frame}


\begin{frame}{Request--Response}
  \tikzstyle{block} = [rectangle, fill=azure(colorwheel), text width=4.5em, text centered, minimum height=4em]
  \tikzstyle{line} = [draw, -latex']
  \tikzstyle{cloud} = [ellipse, fill=wildwatermelon, text width=5em, text centered]
  \tikzstyle{startstop} = [rectangle, rounded corners, text width=5em, minimum height=4em, text centered, fill=tiffanyblue]
  
  \begin{adjustbox}{max totalsize={.9\textwidth}{.8\textheight},center} 
    \begin{tikzpicture}[node distance = 3.5cm]
      \node [cloud] (client1) {Client (Firefox)};
      \node [block, right of=client1] (request) {Request};
      \node [block, right of=client1] (request) {Request};
      \node [cloud, right of=request] (server1) {Server (Apache)};
      \node [block, below of=server1] (generate) {Generate response};
      \node [cloud, below of=generate] (server2) {Server (Apache)};
      \node [block, left of=server2] (response) {Response};
      \node [cloud, left of=response] (client2) {Client (Firefox)};
      % Draw edges
      \path [line] (client1) -- (request);
      \path [line] (request) -- (server1);
      \path [line,dashed] (server1) -- (generate);
      \path [line,dashed] (generate) -- (server2);
      \path [line] (server2) -- (response);
      \path [line] (response) -- (client2);
    \end{tikzpicture}
  \end{adjustbox}
\end{frame}


\begin{frame}{Uniform Resource Locator}
  \textcolor{amethyst}{http}://\textcolor{vermilion}{www}.\textcolor{azure(colorwheel)}{reddit.com}:\textcolor{brightpink}{80}\textcolor{green(ryb)}{/r/funny/}?\textcolor{harvardcrimson}{limit=1}
  
  \begin{table}
    \begin{tabular}{r@{\hspace{0.5cm}}p{6cm}}
      \textcolor{amethyst}{http} & Protocol \\
      \textcolor{brightpink}{80} & Port \\
      \textcolor{vermilion}{www} & Subdomain \\
      \textcolor{azure(colorwheel)}{reddit.com} & Domain \\
      \textcolor{green(ryb)}{/r/funny/} & Path \\
      \textcolor{harvardcrimson}{limit=1} & Parameter
    \end{tabular}
  \end{table}
\end{frame}


\begin{frame}{Resources}
 \begin{columns}[onlytextwidth]
   \begin{column}{0.5\textwidth}
     \centering
      \begin{figure}
        \begin{tikzpicture}
          \draw (0,0) -- (0,1.2) -- (0.7,1.2) -- (0.7,0.8) -- (1,0.8) -- (1,0) -- cycle;
          \draw (0.7,1.2) -- (1,0.8);
        \end{tikzpicture}
        \caption*{File}
      \end{figure}
    \end{column}
    \begin{column}{0.5\textwidth}
      \centering
      \begin{figure}
      \begin{tikzpicture}
        \node (A) [cylinder, shape border rotate=90, draw,minimum height=1.3cm,minimum width=1cm] {};
      \end{tikzpicture}
      \caption*{Database}
      \end{figure}
    \end{column}
  \end{columns}
  
  \begin{quote}
    HTTP is used to transmit resources \ldots A resource is some chunk of information that can be identified by a URL \ldots The most common kind of resource is a file, but a resource may also be a dynamically-generated query result \ldots
  \end{quote}
  \citeurl{www.jmarshall.com/easy/http}
\end{frame}


\begin{frame}{HTTP Methods}
  \begin{description}
    \item[GET] Retrieve information from the server.
    \item[HEAD] Like get, but retrieve only the response header.
    \item[POST] Send data to the server.
    \item[PUT] Set the resource at the URL to the request data.
    \item[DELETE] Delete the resource at the URL.
    \item[CONNECT] Set up tunnel for other traffic to pass through HTTP.
    \item[OPTIONS] Find the allowable operations at the given URL.
    \item[TRACE] Echo the received request.
    \item[PATCH] Partial resource modification.
  \end{description}
\end{frame}

\begin{frame}{Security}
  \begin{itemize}
    \item HTTP is not encrypted.
    \item HTTPS is a protocol based on HTTP, but it provides security.
    \item GET and POST are by far the most commonly used HTTP methods (by web developers).
    \item Data sent by GET and POST will be encrypted over HTTPS.
    \item However, it's generally accepted that POST is more secure for sending sensitive data.
    \item This is because browsers will typically cache and server will typically log URLS, with the data encoded in them.
  \end{itemize}
\end{frame}

\begin{frame}[fragile]{Request and Response Format}
  Requests and responses both have this format:
  \begin{itemize}
    \item Intial line.
    \item Zero or more header lines.
    \item A blank line.
    \item Optional message body (e.g. a HTML file)
  \end{itemize}
  \citeurl{www.jmarshall.com/easy/http}
\end{frame}


\begin{frame}[fragile]{Request (GET) Example}
  \begin{minted}{http}
GET /path/file.html HTTP/1.0
From: someuser@jmarshall.com
User-Agent: HTTPTool/1.0
  \end{minted}
\end{frame}

\begin{frame}[fragile]{Request (POST)}
  \begin{minted}{http}
  POST /path/script.cgi HTTP/1.0
  From: frog@jmarshall.com
  User-Agent: HTTPTool/1.0
  Content-Type: application/x-www-form-urlencoded
  Content-Length: 32
  
  home=Cosby&favorite+flavor=flies
  \end{minted}
\end{frame}

\begin{frame}[fragile]{Response Example}
  \begin{minted}{http}
HTTP/1.1 200 OK
Date: Mon, 27 Jul 2009 12:28:53 GMT
Server: Apache/2.2.14 (Win32)
Last-Modified: Wed, 22 Jul 2009 19:15:56 GMT
Content-Length: 88
Content-Type: text/html
Connection: Closed

<html>
  <body>
    <h1>Hello, World!</h1>
  </body>
</html>
  \end{minted}
\end{frame}



\begin{frame}{URL encoding}
  HTML form data is usually URL-encoded by changing;
  \begin{itemize}
    \item Unsafe characters to \% \emph{xx} where \emph{xx} is the ASCII value.
    \item All spaces to plusses.
    \item Names and values to: name1=value1\&name2=value2.
  \end{itemize}
  
  \begin{description}
    \item[GET] Parameters go in the URL after ?, e.g.\ http://www.google.ie?q=Funny+cats.
    \item[POST] Parameters go in the body.
  \end{description}
  \citeurl{www.jmarshall.com/easy/http}
\end{frame}

\section{REST}

\begin{frame}{REST}
  \begin{itemize}
    \item REST stands for Representational State Transfer.
    \item REST is an architecture describing how we might use HTTP.
    \item RESTful APIs make use of more HTTP methods than just GET and POST.
    \item Most HTTP APIs are not RESTful.
    \item RESTful APIs adhere to a few loosely defined constraints.
    \item Two of those constraints are that the API is stateless and cachable.
  \end{itemize}
  \citeurl{drdobbs.com/web-development/restful-web-services-a-tutorial/240169069}
\end{frame}

\begin{frame}{Stateless}
  \begin{itemize}
    \item Statelessness is a REST constraint.
    \item HTTP uses the client-server model.
    \item The server should treat each request as a single, independent transaction.
    \item No client state should be stored on the server.
    \item Each request must contain all of the information to perform the request. 
  \end{itemize}
\end{frame}

\begin{frame}{Cachable}
  \begin{itemize}
    \item REST APIs should provide responses that are cachable.
    \item Intermediaries between the client and server should be able to cache responses.
    \item This should be transparent to the client.
    \item Cachability increases response time.
    \item Browsers usually cache resources, in case they are requested again.
    \item There is usually a time limit on cached resources.
  \end{itemize}
\end{frame}

\begin{frame}{Example}
  Suppose we have a system for storing and retrieving member's details.
  \begin{table}
    \begin{tabular}{r@{\hspace{0.5cm}}|l@{\hspace{0.5cm}}|l}
      Method & URL & Description \\
      \hline
      GET    & /people   & list members \\
      POST   & /people   & create member \\
      GET    & /people/1 & retrieve member \\
      PUT    & /people/1 & update member \\
      DELETE & /people/1 & delete member
    \end{tabular}
  \end{table}
  \citeurl{bitworking.org/news/201/restify-daytrader}
\end{frame}


\section{JSON}

\begin{frame}{JSON}
  \begin{description}
    \item[JavaScript] A scripting/programming language.
    \vspace{0.25cm}
    \item[Object] Groups of name--value pairs.
    \vspace{0.25cm}
    \item[Notation] Set of rules for representing objects.
  \end{description}
\end{frame}

\begin{frame}{Sending JSON}
  \tikzstyle{block} = [rectangle, fill=azure(colorwheel), text width=4.5em, text centered, minimum height=4em]
  \tikzstyle{line} = [draw, -latex']
  \tikzstyle{cloud} = [ellipse, fill=wildwatermelon, text width=5em, text centered]
  \tikzstyle{startstop} = [rectangle, rounded corners, text width=5em, minimum height=4em, text centered, fill=tiffanyblue]
  
  \begin{adjustbox}{max totalsize={.9\textwidth}{.6\textheight},center} 
    \begin{tikzpicture}[node distance=4cm]
    \node [cloud] (object1) {Object Instance};
    \node [block, below of=object1] (json1) {JSON};
    \node [block, right of=json1, node distance = 6cm] (json2) {JSON};
    \node [cloud, above of=json2] (object2) {Object Instance};
    % Draw edges
    \path [line] (object1) -- node[style={rectangle,fill=white,draw}]{stringify} (json1);
    \path [line] (json1) -- node[style={rectangle,fill=white},draw ]{HTTP} (json2);
    \path [line] (json2) -- node[style={rectangle,fill=white,draw}]{parse} (object2);
    % Draw Memory
    \draw [color=gray, dashed](-2,-1.5) rectangle (2,1.25);
    \node at (-1.4,-1.35) [] {\tiny Machine 1};
    \draw [color=gray, dashed](4,-1.5) rectangle (8,1.25);
    \node at (4.6,-1.35) [] {\tiny Machine 2};
    \end{tikzpicture}
  \end{adjustbox}
\end{frame}

\begin{frame}[fragile]{JSON Example}
  \begin{minted}{json}
{
  "employees": [
    {"firstName":"John", "lastName":"Doe"},
    {"firstName":"Anna", "lastName":"Smith"},
    {"firstName":"Peter", "lastName":"Jones"}
  ]
}
  \end{minted}
\end{frame}

\begin{frame}[fragile]{Using JSON in JavaScript}
  \begin{minted}{javascript}
// Turning text into a JavaScript object.
var obj = JSON.parse(text);
// obj is an obect.

// Turning a JavaScript object into text.
var text = JSON.stringify(obj);
// text is a string.
  \end{minted}
\end{frame}

\begin{frame}{JSON Syntax}
  \begin{itemize}
    \item Name/Value pairs separated by a colon. \\
    \hspace{0.5cm} \mintinline{json}{"name": "Ian"}
    \item Objects identified by curly braces. \\
    \hspace{0.5cm} \mintinline{json}{{}}
    \item Lists identified by square brackets. \\
    \hspace{0.5cm} \mintinline{json}{[]}
    \item All strings (and names) use double quotes (not single). \\
    \hspace{0.5cm} \mintinline{json}{"Ian"}
  \end{itemize}
\end{frame}

\begin{frame}{JSON Types}
  \begin{itemize}
    \item Numbers \\
    \hspace{0.5cm} \mintinline{json}{123.456}
    \item Strings \\
    \hspace{0.5cm} \mintinline{json}{"Hello, world!"}
    \item Boolean \\
    \hspace{0.5cm} \mintinline{json}{true"}
    \item Arrays\\
    \hspace{0.5cm} \mintinline{json}{[1,2,3]}
    \item Objects\\
    \hspace{0.5cm} \mintinline{json}{{"name": "Ian"}}
    \item null \\
    \hspace{0.5cm} \mintinline{json}{null}
  \end{itemize}
\end{frame}

\section{XML}

\section{AJAX}

\section{NoSQL}

\section{Map Reduce}